\chapter{Statistical Inferences for Proportions}

\section{Interval Estimation and the Confidence Interval}

Suppose you had to guess your own weight, and then guess a classmate's weight. You would probably feel more confident about your own guess. You might say 140 pounds for both, but you would trust your own number a lot more than the one for your classmate. A point estimate only gives you a single number. It does not tell you how much to trust that number.

A point estimate is a single value that is calculated from a sample. It serves as our best guess for a population parameter. Although it can be useful, it is incomplete. Two point estimates can look identical on paper while carrying very different amounts of certainty behind them like the way your weight guess and your classmate's guess might both be 140 pounds despite being backed by very different amounts of information.

This is where the \newterm{confidence interval}\index{confidence interval} comes in. A confidence interval gives a range of plausible values for the population parameter that is built around the point estimate. The \newterm{margin of error}\index{margin of error} is the "plus or minus" amount that sets the width of that range. It is driven by how confident we want to be and how much the estimate varies from sample to sample.

% Think about this: if someone tells you "I am 95\% confident the true proportion is between 0.60 and 0.70," what do you think happens to that range if they instead wanted to be 99\% confident? Do you think the interval gets wider or narrower?

A wider interval is more likely to capture the true value, but it also tells you less. This tradeoff between confidence and precision runs through everything in this chapter.

% PLACEHOLDER: Conditions for Inference, random sampling, 10% condition (independence), and  large counts condition (n\hat{p} >= 10 and n(1-\hat{p}) >= 10). Add notes/drafts here.


\section{Constructing a Confidence Interval for a Proportion}

When we discussed confidence intervals, we saw that a wider interval gives us more certainty at the cost of being precise. That tradeoff is captured by a single number called the \newterm{critical value}\index{critical value}, written $z^*$.

Let’s say a city council wants to estimate the proportion of residents who support building a new network of bike lanes. Before collecting any data, they need to decide how confident they want to be in their eventual estimate, say 95\%. That confidence level determines $z^*$.

A 95\% confidence level leaves 5\% split evenly between the two tails of the standard normal distribution, so 2.5\% in each tail. We need the value $z^*$ such that the area to its right is 0.025, which means the area to its left is 0.975. Using a standard normal table or a calculator's inverse normal function, $z^* \approx 1.96$.

If the city council instead wanted 99\% confidence, would $z^*$ be larger or smaller than 1.96?

Once we have $z^*$, we still need to decide how precise we want our estimate to be. An interval can technically be true and still be worthless. If the council's confidence interval for bike lane support turned out to be "somewhere between 5\% and 95\%," that would be almost useless for making a decision, even if it is a perfectly valid 95\% confidence interval. The width of a useful interval, called the margin of error, depends on how large a sample the council collects. Larger samples produce narrower, more informative intervals for the same confidence level.

% [PLACEHOLDER: General confidence interval formula for a proportion.

% [PLACEHOLDER: Sample size formula (solving for n).

This is why sample size planning matters just as much as the confidence level. A researcher who wants a useful interval has to think ahead about how much data collecting that precision will require.

\textbf{Computing in Excel}

% [Excel walkthrough placeholder: compute z* using NORM.S.INV(0.975),
%   construct the confidence interval using cell formulas for p-hat,
%   the standard error, and the margin of error, and compute required
%   sample size using the solved-for-n formula above. Add screenshots
%   and step-by-step instructions here once ready. ]

\textbf{Computing in Python}

% [PLACEHOLDER: Python walkthrough using scipy.stats.norm.ppf() to find
%   z*, statsmodels.stats.proportion.proportion_confint() to construct
%   the confidence interval directly from counts, and a short function or
%   manual calculation to solve for required sample size given a target
%   margin of error. Add code, output, and interpretation here once
%   notes/drafts are available. ]


\section{Inference: Tests of Significance}

Building a confidence interval answers "what is the population proportion likely to be?" But sometimes the question you actually care about is different: has something changed? \newterm{Statistical inference}\index{Statistical inference} is the process of using information from a sample to answer a specific question or make a decision about a population parameter, and a \newterm{test of significance}\index{test of significance} is the tool built for exactly this kind of question.
